\documentclass[article,12pt]{elsarticle}
\usepackage[spanish]{babel}
\usepackage[utf8]{inputenc}
\usepackage{amssymb}
\usepackage{flushend}
\usepackage{float}
\usepackage{anysize}
\usepackage{amsmath}
\usepackage{amsfonts}
\usepackage{dcolumn}
\usepackage{amsthm}
\usepackage{caption}
\usepackage{float}
\usepackage{subfigure}
\usepackage{graphicx}
\usepackage{lmodern}
\usepackage{booktabs}
\usepackage{color}
\usepackage{longtable}
\usepackage{rotating}
\usepackage{etoolbox}
\usepackage{comment}
\usepackage{pdflscape}
\usepackage{booktabs}


\begin{document}

\begin{frontmatter}
    \title{Laboratorio de Fenómenos Colectivos \\ Temperatura del Fuego} \\
    \author{Andrés López González$^1$} %Aquí va el nombre de los integrantes que participaron en el informe.
    \address{$^1$ Facultad de Ciencias, Universidad Nacional Autónoma de México, M\'exico CDMX.}
\end{frontmatter}

Dado que esta actividad solo pide resultados iré directo a estos y sus interpretaciones; primero nos interesa obtener el calor específico del balín, para ello obtenemos la siguiente tabla.

\begin{table}[h]
\centering
\begin{tabular}{|l|l|}
\hline
Temperatura agua al tiempo & $(19 \pm 0.5)^{\circ}$ \\ \hline
Temperatura balín frío     & $(-6 \pm 0.5)^{\circ}$ \\ \hline
Temperatura de equilibrio  & $(18 \pm 0.5)^{\circ}$ \\ \hline
Masa del agua              & $(24.17 \pm 0.005)g$   \\ \hline
Masa del balín             & $(8.39 \pm 0.005)g$    \\ \hline
\end{tabular}
\caption{Primera fase}
\end{table}

De aquí tenemos lo siguiente usando la ecuación de calor cedido-obtenido.

\begin{align*}
    m_bc_b(T_{eq}-T_f) &= m_ac_a(T_{amb}-T_{eq}) \\
    c_b &= \frac{m_ac_a(T_{amb}-T_eq)}{c_b(T_{eq}-T_f)} \\
    &= \frac{(24.17)(4.18)(19 - 18)}{(8.39)(18 - (-6))} \\
    &= 0.501 \, \text{J/g°C}
\end{align*}

Para la segunda fase obtuvimos esta siguiente tabla, en donde dejamos al balín en fuego directo durante más de 15 minutos

\begin{table}[h]
\centering
\begin{tabular}{|l|l|}
\hline
Temperatura agua al tiempo & $(19 \pm 0.5)^{\circ}$ \\ \hline
Temperatura del fuego     & $(? \pm 0.5)^{\circ}$ \\ \hline
Temperatura de equilibrio  & $(29 \pm 0.5)^{\circ}$ \\ \hline
Masa del agua              & $(45.18 \pm 0.005)g$   \\ \hline
Masa del balín             & $(8.39 \pm 0.005)g$    \\ \hline
\end{tabular}
\caption{Segunda fase}
\end{table}

Usamos la misma ecuación del ejercicio anterior, aquí nuestra f en vez de frío es de fuego.

\begin{align*}
    m_bc_b(T_{eq}-T_f) &= m_ac_a(T_{amb}-T_{eq}) \\
    T_f &= T_{eq}-\frac{m_ac_a(T_{amb}-T_{eq})}{m_bc_b} \\
    &= 29 - \frac{(45.18)(4.18)(19 - 29)}{(8.39)(0.501)} \\
    &= 478.4 \, \text{°C}.
\end{align*}


\end{document}

